% Ilustración vectorial original: esquema simplificado, no a escala.
% Se inserta dentro de un frame con \input; usa la paleta de la presentación.
\begin{tikzpicture}[x=0.89cm,y=0.89cm,>=Latex,
    every node/.style={font=\scriptsize},
    dendrita/.style={draw=verde,line width=1.5pt,line cap=round},
    guia/.style={draw=fondo!65,line width=0.5pt}]
    \path[use as bounding box] (-0.1,-1.95) rectangle (8.2,2.15);

    % Arborización dendrítica: recibe conexiones de otras neuronas.
    \draw[dendrita] (2.35,0.3) .. controls (1.85,0.7) and (1.55,0.8) .. (1.15,1.15);
    \draw[dendrita] (1.65,0.8) -- (1.45,1.45) -- (1.1,1.7);
    \draw[dendrita] (1.45,1.45) -- (1.65,1.8);
    \draw[dendrita] (1.15,1.15) -- (0.6,1.3) -- (0.3,1.65);
    \draw[dendrita] (0.6,1.3) -- (0.2,1.2);
    \draw[dendrita] (2.15,0.05) .. controls (1.5,0.25) and (1.2,0.1) .. (0.7,0.3);
    \draw[dendrita] (0.7,0.3) -- (0.25,0.6);
    \draw[dendrita] (0.7,0.3) -- (0.1,0.15);
    \draw[dendrita] (2.2,-0.25) .. controls (1.6,-0.6) and (1.35,-0.5) .. (0.85,-0.95);
    \draw[dendrita] (1.4,-0.6) -- (0.75,-0.5) -- (0.35,-0.7);
    \draw[dendrita] (0.85,-0.95) -- (0.65,-1.35);
    \draw[dendrita] (0.85,-0.95) -- (0.3,-1.1);
    \draw[dendrita] (2.45,-0.4) -- (2.05,-1.05) -- (1.6,-1.35);
    \draw[dendrita] (2.05,-1.05) -- (2.15,-1.45);

    % Soma y núcleo. La forma irregular evoca una neurona multipolar.
    \path[draw=fondo,fill=fondo!12,line width=1.1pt]
        (2.1,0.05)
        .. controls (2.25,0.2) and (2.18,0.55) .. (2.4,0.65)
        .. controls (2.7,0.5) and (2.95,0.65) .. (3.12,0.35)
        .. controls (3.28,0.1) and (3.3,0.05) .. (3.6,0)
        .. controls (3.3,-0.05) and (3.25,-0.3) .. (3.05,-0.5)
        .. controls (2.85,-0.65) and (2.65,-0.5) .. (2.43,-0.68)
        .. controls (2.45,-0.35) and (2.18,-0.3) .. (2.1,0.05) -- cycle;
    \filldraw[fill=white,draw=fondo!60] (2.72,-0.02) circle (0.24);
    \fill[fondo!55] (2.75,0.02) circle (0.07);

    % Axón con segmentos de mielina y arborización terminal.
    \draw[draw=fondo,line width=1.8pt] (3.5,0) -- (6.65,0);
    \foreach \inicio in {3.7,4.5,5.3} {
        \draw[draw=fondo!65,fill=fondo!12,rounded corners=3pt,line width=0.8pt]
            (\inicio,-0.16) rectangle ++(0.6,0.32);
    }
    \draw[draw=naranja,line width=1.5pt,line cap=round]
        (6.65,0) -- (7.05,0.7) -- (7.65,1.0)
        (7.05,0.7) -- (7.3,1.35)
        (6.65,0) -- (7.15,0.05) -- (7.8,0.3)
        (7.15,0.05) -- (7.75,-0.2)
        (6.65,0) -- (7.05,-0.65) -- (7.65,-1.05)
        (7.05,-0.65) -- (7.25,-1.35);
    \foreach \punto in {(7.65,1.0),(7.3,1.35),(7.8,0.3),(7.75,-0.2),(7.65,-1.05),(7.25,-1.35)} {
        \fill[naranja] \punto circle (0.075);
    }

    % Etiquetas anatómicas y dirección conceptual de la señal.
    \node[text=verde,font=\scriptsize\bfseries] at (0.7,2.0) {Dendritas};
    \draw[guia] (0.7,1.82) -- (0.95,1.42);
    \node[align=center,font=\scriptsize\bfseries] at (3.1,1.55) {Cuerpo celular};
    \draw[guia] (3.1,1.35) -- (2.95,0.63);
    \node[font=\scriptsize\bfseries] at (4.95,0.85) {Axón};
    \draw[guia] (4.95,0.65) -- (4.95,0.2);
    \node[text=naranja,align=center,font=\scriptsize\bfseries] at (7.4,1.95) {Terminales};
    \draw[guia] (7.4,1.75) -- (7.35,1.43);
    \draw[->,fondo!65,line width=0.8pt] (2.65,-1.7) -- (6.65,-1.7)
        node[midway,below=1mm] {sentido de propagación};
\end{tikzpicture}
